%! Author = sriadityavedantam
%! Date = 3/25/26

\lesson{12}{Wednesday 10 September 2025 9:10}{Day 12}

\begin{theorem}[Bolzano-Weierstrass]{
    If $(x_n)$ is bounded, then there exists a subsequence converging to $\limsup x_n = L$.
}
    To construct a sequence $(m_k)$ in $\n$ such that $m_k\to\infty$ and $x_{m_k}\to L$, let
    \[
        y_k \coloneqq \sup\{x_m : m\geq k\}.
    \]
    Then $y_k\to L$, by definition.
    For each $k$, choose $m_k\geq k$ such that
    \[
        y_k\geq x_{m_k}\geq y_k - \dfrac{1}{k}.
    \]
    Since $y_k\to L$, then
    \begin{align*}
        \lim\left(y_k - \dfrac{1}{k}\right) &= \lim y_k - \lim \dfrac{1}{k}\\
        &= L - 0\\
        &= L.
    \end{align*}
    Hence, by the Squeeze Theorem, $x_{m_k}\to L$.
\end{theorem}

\section{The Cauchy Criterion}

\begin{definition}[Cauchy Sequences]
    A sequence $(x_n)$ in $\r$ is Cauchy if for all $\eps > 0$, there exists $N\in\n$ such that if $m,n\geq N$, then $\abs{x_m - x_n} < \eps$.
\end{definition}

\begin{proposition}{
    Any convergent sequence is Cauchy.
}
    Let $\eps > 0$ and choose $N\in\n$ such that:
    \begin{align*}
        n > N&\implies \abs{x_n - L} < \dfrac{\eps}{2}\\
        m > N&\implies \abs{x_m - L} < \dfrac{\eps}{2}.
    \end{align*}
    Then $m,n > N$ implies:
    \begin{align*}
        \abs{x_m - x_n} &= \abs{x_m - L + L - x_n}\\
        &\leq \abs{x_m - L} + \abs{x_n - L}\\
        &< \dfrac{\eps}{2} + \dfrac{\eps}{2} = \eps.
    \end{align*}
\end{proposition}

\begin{lemma}{
    Any Cauchy sequence is bounded.
}
    Let $N\in\n$ such that if $m,n \geq N$, then $\abs{x_m - x_n}<1$.
    Then for $m\geq N$,
    \[
        \abs{x_m} = \abs{x_m - x_N + x_N} \leq \abs{x_m - x_N} + \abs{x_N} < 1 + \abs{x_N}.
    \]
    Hence, for any $m\in\n$, we have
    \[
        \abs{x_m} \leq\max\{\abs{x_N} + 1, \abs{x_1},\abs{x_2}, \ldots, \abs{x_N}\}.
    \]
\end{lemma}

\begin{theorem}{
    Any Cauchy sequence converges.
}
    Define
    \[
        y_n \coloneqq \sup\{x_m : m\geq n\}
        \quad\text{and}\quad
        z_n \coloneqq \inf\{x_m : m\geq n\}.
    \]
    Let $\eps > 0$.
    Choose $N\in\n$ such that if $m,n\geq N$, then $\abs{x_n - x_m} < \eps$.
    Set $S_N \coloneqq \{x_n : n\geq N\}$.
    Then for all $a,b\in S_N$, we have
    \[
        \abs{a-b}<\eps.
    \]
    Hence
    \[
        y_N - z_N = \sup(S_N) - \inf(S_N) \leq \eps.
    \]
    Therefore, for all $n\geq N$,
    \[
        0\leq y_n - z_n \leq y_N - z_N \leq \eps.
    \]
    So $\lim (y_n - z_n) = 0$.
    Since $z_n \leq x_n \leq y_n$ for all $n$, and since $(y_n)$ is decreasing while $(z_n)$ is increasing, both are monotone and bounded, hence convergent.
    Let
    \[
        \lim y_n = A
        \quad\text{and}\quad
        \lim z_n = B.
    \]
    Since $\lim (y_n-z_n)=0$, we have $A-B=0$, so $A=B$.
    Hence, by the Squeeze Theorem, $(x_n)$ converges.
\end{theorem}